\documentclass[12pt, a4paper]{article}

\usepackage{setspace} % Necessary for setting space
\usepackage{amsmath}
\usepackage{amssymb}
\usepackage{harvard} % Necessary for Econometrica style reference
\usepackage{hyperref} % Necessary for citing websites
\usepackage{lasstyle}
\usepackage{booktabs} % Table
\usepackage{graphicx} % Table
\usepackage{amsthm}

\newtheorem{proposition}{Proposition}

\newcounter{headercounter}
\setcounter{headercounter}{1}

\newcommand{\header}[1]{\begin{center}{\large \textbf{\Roman{headercounter}. #1}} \end{center} \stepcounter{headercounter}}

% \renewcommand{\thesection}{\Roman{section}} % Roman section enumeration

\doublespacing

\author{Lucas Sernik}
\title{The implications of spending pattern differences between mayors who might be reelected and those who are bound by term limits: A case study from municipalities in São Paulo}
\date{\today}

\begin{document}
\maketitle

\begin{abstract}
This paper examines whether reelection incentives distort the composition of municipal spending in São Paulo, Brazil. Using a panel of 643 municipalities over 2017--2024 and exploiting variation in incumbent reelection eligibility generated by Brazil's two-term limit rule, I estimate a time fixed effects model relating reelection eligibility to the share of expenditure allocated to six functional categories. I find that mayors who can run for reelection spend more on administration and less on urban infrastructure. These findings are consistent with rent extraction on urban projects when mayors cannot be reelected, and with political patronage when they can still run again. Also, mayors who can be reelected and have below-average discretion over budget spending spend more on administration and social assistance, suggesting that budget is allocated to these categories to gain advantages in the next election.
\end{abstract}

\newpage

\header{Introduction}

In a democratic society, citizens vote for a set of candidates for government positions who are supposed to make decisions that promote the public good. Recent developments in Political Science and Political Economy offer different views on whether the existence of term limits provides any societal benefit. Smart and Sturn (2013) develop a game that models voters and self-interested politicians, in which term limits help voters screen out politicians who use their position for their own self-interest, improving overall welfare despite allowing sub-optimal behavior on the binding term. Fowler (2024) discusses how term limits might have unintended effects when applied to political positions. He argues that term limits might remove from office a high-quality official that voters like and who voters would prefer to stay in office, and, more importantly, that the mere existence of a term limit modifies the incentives that the official has to do their job. Such concern is tested by Besley and Case (1995). 

Their study that provides empirical evidence to the vast theoretical literature on politicians and what drives their decisions. They argue that the existence of term limits might distort politicians' incentives concerning which policies are implemented. After grounding their theoretical framework and some of its shortcomings, they analyze state governors from 1950 to 1986 and conclude that ``lame ducks'', governors who cannot run for the next term, increase their spending and taxation in their last term, since they generally do not need to maintain their political reputation and are prone to making unpopular decisions. They point out that term limits might have adverse effects on economic policy, arguing that politicians might behave differently and be more inefficient due to the different incentives that such a rule imposes on their constrained maximization problem.

This paper will focus on Brazil, more specifically, on municipalities in the state of São Paulo. I will test whether the presence of the binding two consecutive term limit on any elected official for a mayor in the state of São Paulo has any impact on how such mayors choose to spend their budget. Brazil transitioned from military rule to a democratic regime at the end of the 1980's, with a movement for direct elections in 1985 and the creation of a new constitution in 1988. Recent efforts for greater transparency in government actions and budgeting led the Brazilian Congress to pass the ``Lei da Transparência'' in 2009 and the ``Lei de Acesso à Informação'' in 2011, which obliges public institutions to publish detailed documents with information on government revenues and expenses on the internet. In the state of São Paulo, expenditure data are available dating back to 2014. This paper will attempt to exploit expense data made available by these laws in the 645 municipalities of the state of São Paulo from 2017 to 2024. This time period is the largest one that contains entire election terms (two) up to 2026. I will employ similar methods to Besley and Case (1995) to test whether mayors bound by term limits behave differently than those who may be reelected in the next period, according to the type of expense and the year. I find that reelection eligibility significantly affects the composition of discretionary spending in two categories. Reelection-eligible mayors allocate approximately 0.94 percentage points less of their discretionary budget to urbanismo relative to term-limited mayors ($p = 0.002$), consistent with electoral accountability constraining rent extraction through inflated infrastructure contracts. They also allocate approximately 0.89 percentage points more to administration ($p = 0.016$), consistent with political patronage. This supports the interpretation that incumbents seeking reelection expand the municipal payroll to generate electoral loyalty. I also find evidence of conditional clientelism: reelection-eligible mayors with below-average fiscal discretion spend more on both administration and social assistance, suggesting that fiscal constraint shapes the instruments of electoral manipulation.

\header{Theoretical Framework and Related Literature}

\subsection*{Political Budget Cycle}

The political budget cycle literature examines how electoral incentives distort fiscal policy, building on the theoretical foundations laid by \citeasnoun{rogoff:1990}, who shows that incumbents facing reelection have incentives to signal competence by manipulating observable fiscal instruments before elections. In Rogoff's model, voters cannot directly observe the government's type, so incumbents exploit this information asymmetry by shifting spending toward more visible categories and away from productive investment in the pre-electoral period. The empirical literature has confirmed the existence of political budget cycles in different institutional contexts. \citeasnoun{shi:2006} provides cross-country evidence that electoral cycles in fiscal balances are larger in countries with less informed voters, which is consistent with the information asymmetry mechanism exploited by politicians in the theoretical model. \citeasnoun{veiga:2007} find evidence of pre-electoral manipulation of local government expenditures in Portuguese municipalities. \citeasnoun{brender:2005} find that political budget cycles are more pronounced in new democracies where voters have less experience evaluating incumbent performance.

More recently, the literature has moved beyond total fiscal expenditure to examine whether electoral incentives distort the composition of spending. \citeasnoun{drazen:2010} argue that the composition channel is theoretically more important than the level channel, since voters may be more sensitive to the type of spending instead of its total amount. \citeasnoun{drazen:2010} also shows empirical evidence of this through analyzing municipalities in Colombia. My paper contributes to this literature by providing direct evidence on how reelection incentives affect discretionary spending composition at the municipal level in São Paulo, exploiting the variation in electoral incentives generated by Brazil's two consecutive term limit rule.

\subsection*{Electoral Accountability and Political Agency}

The electoral accountability literature asks whether the threat of electoral punishment disciplines incumbent politicians and induces them to act in voters' interests. The theoretical framework builds on principal-agent models of politics in which voters --- the principals --- delegate policy decisions to politicians --- the agents --- but cannot perfectly observe their effort or type. \citeasnoun{barro:1973} and \citeasnoun{ferejohn:1986} establish the foundational theoretical result that reelection incentives can discipline politicians: incumbents who care about remaining in office will exert effort and restrain opportunistic behavior to avoid electoral punishment.

The core empirical contribution in this literature is \citeasnoun{besley:1995}, who exploit variation in gubernatorial term limits across U.S. states to test whether the absence of reelection incentives, made possible by binding term limits, alters economic policy choices. They find that term-limited governors raise taxes and increase spending, consistent with a model in which reelection-eligible incumbents exert effort to keep taxes and spending low to signal competence. My paper adopts the same identification strategy, using Brazil's two-term mayoral limit to generate exogenous variation in reelection eligibility, and extends the Besley and Case framework in two directions. First, we focus on spending composition rather than aggregate levels, asking not whether mayors spend more or less but which categories receive more funds from the discretionary budget. Second, I explicitly account for the degree of fiscal discretion available to each mayor as a moderating factor, recognizing that the scope for politically motivated spending decisions depends on how much of the budget is subject to mayoral choice rather than mandated earmarking, which is common for many municipalities in Brazil.

Subsequent work in this tradition has extended the Besley and Case framework to a variety of settings. \citeasnoun{ferraz:2011} exploit randomized audits of Brazilian municipalities to provide direct evidence that electoral accountability reduces corruption, establishing that the threat of electoral punishment disciplines not just aggregate fiscal behavior but also the integrity of public procurement. \citeasnoun{brollo:2013} show that larger federal transfers to Brazilian municipalities increase corruption and reduce the quality of politicians, highlighting how fiscal capacity interacts with political incentives. \citeasnoun{brollo:2012} show that federal transfers to Brazilian municipalities are systematically higher when the mayor is politically aligned with the federal government, which demonstrates how fiscal transfers in Brazil are subject to political manipulation beyond the electoral cycle mechanism. My contribution to this literature is to show that electoral accountability operates through specific spending categories in the discretionary budget, and that the form of electoral manipulation depends on the fiscal constraints facing each mayor and their political intentions, a mechanism I term conditional clientelism.

The reputation-building model of political behavior presented in Besley and Case (1995) provides a theoretical foundation for understanding how electoral incentives shape incumbent behavior, which will also be used in this paper. While Besley and Case (1995) applied it to governors, I apply it to mayors.

Each incumbent mayor is characterized by an unobservable type $\omega \in \{\omega_1, ..., \omega_N\}$, ordered $\omega_1 < ... < \omega_N$, which captures their intrinsic competence or commitment to public service. The probability that a mayor is of type $\omega_j$ is denoted $\pi_j$. While in office, the mayor takes an unobservable action $\alpha \in [\underline{\alpha}, \bar{\alpha}]$, interpreted as the amount of effort invested in effective governance. This effort probabilistically affects a policy outcome $r \in \mathbb{R}^+$ that voters care about, distributed according to $F(r \; | \; \alpha)$. Higher effort makes better $r$ outcomes more likely.

The mayor derives utility $v(\alpha, \omega)$ from holding office and exerting effort $\alpha$. We assume that $v(\alpha, \omega)$ is decreasing and concave in $\alpha$, increasing in $\omega$, and that the cross-derivative between $\alpha$ and $\omega$ is positive, so that higher-type mayors derive more utility from exerting effort and acting in the public interest. The mayor receives zero utility when out of office.

The model compares mayoral behavior under two regimes. In the term-limited regime, the mayor cannot run for reelection and therefore faces no future electoral consequences. The optimal action simply maximizes immediate payoffs, represented by $\alpha_s(\omega)$:

\begin{equation}
    \alpha_s(\omega) = \arg\max_{\alpha \in [\underline{\alpha}, \bar{\alpha}]} \left\{ v(\alpha, \omega) \right\}
\end{equation}

This is the \textit{short-run} decision ($s$ subscript). Since effort is costly and there is no reputational benefit to exerting more effort than is immediately optimal, term-limited mayors choose the action that maximizes current, private utility alone.

In the reelection-eligible regime, the mayor cares about the future and their reelection probability. Voters observe the policy outcome $r$ and use it to update their beliefs about the mayor's type via Bayes' rule.  The probability of such policy outcome $r$ happening depends on the amount of effort $\alpha_\ell(\omega_k)$ that a mayor exerted given their type $\omega_k$, specifically considering that they can run for reelection (subscript $\ell$). Mathematically, the posterior probability that the mayor is of type $\omega_k$ given outcome $r$ is:

\begin{equation}
    \beta_k(r) = \frac{\pi_k f(r \; | \; \alpha_l(\omega_k))}{\sum_{j=1}^{N} \pi_j f(r \; | \; \alpha_l(\omega_j))}
\end{equation}

Voters use a cutoff rule: there exists a threshold $r^*$ such that the mayor is reelected if and only if $r \geq r^*$. The reelection-eligible mayor therefore solves:

\begin{equation}
    \alpha_l(\omega) = \arg\max_{\alpha \in [\underline{\alpha}, \bar{\alpha}]} \left\{ v(\alpha, \omega) + \delta \Pr\left[r \in R(\sigma) \mid \alpha\right] v(\alpha_s(\omega), \omega) \right\}
\end{equation}

where $\delta$ is the discount factor, $R(\sigma) = \{r : \sigma(r) = 1\}$ is the set of outcomes that allow reelection under optimal voters' strategy $\sigma$, and $\Pr[r \in R(\sigma) \mid \alpha]$ is the probability of reelection given effort $\alpha$. The second term represents the expected future payoff from winning reelection, discounted to the present through $\delta$. Since higher effort increases the probability that $r \geq r^*$, reelection-eligible mayors have an incentive to exert more effort than their term-limited counterparts and please the voters.

The central prediction of the model is:

\begin{proposition}
If two terms are allowed, incumbents who are eligible for reelection exert more effort than term-limited incumbents. Those in their last term give lower payoffs to voters, on average, relative to their first term in office.
\end{proposition}

Reelection-eligible mayors internalize the reputational return to effort. By delivering better outcomes today, they signal to voters that they are a high type and increase their probability of securing a second term. Term-limited mayors face no such incentive and revert to the short-run optimum, which involves less effort and consequently lower expected payoffs to voters.

\subsection*{Implications for spending composition}

While the original model is cast in terms of aggregate effort and outcomes, its predictions extend to the composition of public spending. A mayor who wishes to signal competence to voters faces incentives to allocate discretionary resources toward categories that are visible, attributable, and electorally salient. Conversely, a term-limited mayor with no future electoral concerns may exploit the discretionary budget for personal or political gain without fear of electoral punishment.

This generates two testable predictions that motivate my empirical specification. First, reelection-eligible mayors should allocate their discretionary budget differently from term-limited mayors, with the direction of the effect depending on which categories are most electorally salient or most vulnerable to rent extraction. Second, the magnitude of this effect should be increasing in the degree of fiscal discretion available to the mayor: when the budget is largely earmarked by constitutional mandates, the scope for politically motivated reallocation is restricted regardless of electoral incentives. I test both predictions using the panel of São Paulo municipalities described in section (IV).

\subsection*{Applicability to São Paulo municipal elections}

The theoretical framework developed by \citeasnoun{besley:1995} was originally applied to U.S. governors, but its core assumptions map naturally onto the Brazilian municipal context, and in particular onto the São Paulo municipalities. I discuss each core assumption.

\textbf{Unobservable effort and observable outcomes.} The model requires that mayoral effort is unobservable to voters, but that policy outcomes, which depend on effort, are at least partially observable. This assumption is well-suited to the Brazilian municipal context. Voters cannot directly observe how a mayor negotiates contracts, manages the municipal bureaucracy, or know their true intentions once they hold office. However, they can observe tangible outcomes such as the quality of local public services, the condition of roads and public works, and the size of the municipal payroll. The information asymmetry that motivates the reputation-building mechanism is therefore present.

\textbf{Binding term limits as an exogenous source of variation.} The key identification assumption is that term limits generate exogenous variation in the discount rate of political agents: mayors in their second and final term no longer face future electoral consequences and therefore have little incentive to exert effort. Brazilian constitutional law limits mayors to two consecutive terms in office, creating a discontinuity in electoral incentives between first-term and second-term mayors within the same election cycle.

\textbf{Fiscal discretion as a moderating condition.} The original model abstracts from the institutional constraints on policy choice, treating effort as freely translatable into outcomes. In the Brazilian municipal context, however, a substantial share of revenue is constitutionally earmarked. Federal transfers through FUNDEB and SUS are tied to education and health, respectively, for example, leaving mayors with limited discretion over a significant portion of their budget. I extend the model's implicit assumption of unconstrained discretion by explicitly conditioning on the share of discretionary spending available to each mayor. This extension is theoretically motivated, since mayors can only change spending patterns if and only if they can influence such spending to some degree. Municipalities where earmarked transfers dominate the budget offer less scope for politically motivated reallocation, regardless of electoral incentives.

\textbf{Heterogeneity in the form of opportunistic behavior.} The model predicts that term-limited mayors revert to short-run payoff maximization, but nothing is said on the specific form this takes. In the Brazilian context, two channels are particularly relevant. First, public works and infrastructure spending have been extensively documented as a vehicle for rent extraction through inflated contracts and kickbacks in Brazilian municipalities \cite{ferraz:2011}. These are captured by expenses in the urbanism category in the dataset. Term-limited mayors, facing no future electoral punishment, may exploit this channel freely. Second, patronage through administrative hiring is a well-documented instrument of electoral clientelism in Brazil. Political connections are one of the major reasons why public employees were selected in Brazil during 1997--2014, primarily due to patronage and leading to the selection of less qualified individuals for a given role \cite{colonnelli:2020}. Reelection-eligible mayors have stronger incentives to invest in this form of patronage and, as my results suggest, may be used as a tool for securing electoral support. Both predictions are testable with my data and constitute the primary empirical contribution of this paper.

\header{Data}

Expense data was collected through the publicly available API from Tribunal de Contas do Estado de São Paulo \cite{TCESP:expenses}. Expense data was categorized according to the official categories determined by São Paulo's Plano Plurianual (PPA), which is the budget plane made by the state of São Paulo to its municipalities. It establishes targets and laws that municipalities within the state must abide. Even though several such categories were analyzed, only a few of them presented meaningful results, which are highlighted in this paper. These categories are Health, Education, Urbanism, Social Assistance, Administration, and Transport. Out of the 645 municipalities in the state of São Paulo, the public API offers expense data for 644, all municipalities except the São Paulo municipality. The São Paulo municipality has their own website and data available from 2021 onwards, reason why it will be dropped from the dataset. Out of the 644 remaining municipalities in the state of São Paulo, the public API offers expense data from 2014 until 2026. Municipal elections in Brazil have happened every four years since 1996 \cite{TSE:eleicao}. This means that municipal elections happened in 2012, 2016, 2020, and 2024. I restricted calls to the years 2017 to 2024, the two complete election periods with relevant data. After downloading with the API, the municipality of Florínea was dropped, because it did not contain expenses for the entire time period. This left us with 643 municipalities, which is the number of municipalities in the analysis.

Election data was collected through the publicly available data from Tribunal Superior Eleitoral \cite{TSE:data}. Through their data transparency tool, I was able to collect mayoral election data such as number of candidates, whether the incumbent may run to reelection or not, and reelection status, which states whether the winner was in their reelection campaign.

Control data was collected from 2 sources. The first is from Instituto Brasileiro de Geografia e Estatística \cite{IBGE:sidra}. The second is from Atlas Brasil \cite{atlas:idh}. I collected population, GDP per capita, basic sanitation services, and age distribution at the municipal granularity from the year 2010, the most recent data available. The second is from Atlas Brasil \cite{atlas:idh}. Descriptive statistics of the data are available in the Appendix.

\header{Estimation and Methods}

\subsection*{Ideal model}

Let $s_{ikt}$ denote the share of total expenditure allocated to category $k$ by municipality $i$ in year $t$. The baseline model is:

\begin{equation*}
    s^{d}_{ikt} = \beta_0
                + \beta_1 R_{it}
                + \beta_2 D_{it}
                + \beta_3 (R_{it} \times D_{it})
                + \beta_4 N_{it}
                + \beta_5 \ln(P_i)
                + \mathbf{X}_i'\boldsymbol{\gamma}
                + \mu_t
                + \varepsilon_{it}
\end{equation*}

\noindent where $s^{d}_{ikt}$ is the discretionary spending share in category $k$, municipality $i$ at year $t$, $R_{it}$ is reelection eligibility for an incumbent in municipality $i$ at year $t$, $D_{it}$ is ratio between discretionary budget and total budget (discretionary budget + earmarked budget), $R_{it} \times D_{it}$ is the interaction term, which captures the effect of changing $D_{it}$ when an incumbent is eligible for reelection, $N_{it}$ is the number of candidates in the relevant election, $\mathbf{X}_i$ is a vector containing 2010 Census controls (PIB per capita, IDHM, share of children), and $\mu_t$ are the year fixed effects.

The coefficient of interest are $\beta_1$, $\beta_2$ and $\beta_3$. $\beta_1$, identifies the effect of reelection eligibility on spending composition. $\beta_2$ explains how much total discretion. The interaction coefficient $\beta_3$ tests whether this effect is moderated by fiscal discretion.

\subsection*{Robustness checks and updated model}

\subsubsection*{Justification for demeaning the discretionary share}

In the updated model, the discretionary share of total expenditure $D_{it}$ enters the regression in demeaned form, $\tilde{D}_{it} = D_{it} - \bar{D}$, where $\bar{D} = 71.4\%$ is the sample mean. This transformation is motivated by the absence of any municipalities where the mayors have no discretion over public funding allocation, represented by $D_{it} = 0$, in the dataset and in the real world. It simply recenters the discretionary share term around the average so that $\beta_1$ represents the reelection effect at the average level of fiscal discretion in the sample rather than at the economically meaningless value of $D_{it} = 0$.

\subsubsection*{Fiscal discretion as a proxy for municipal wealth}

A potential concern with my measure of fiscal discretion is that $D_{it}$ may simply proxy for municipal wealth. Richer municipalities tend to generate more own-source tax revenue, which increases the discretionary share of their budget. If this were the case, the interaction term $R_{it} \times \tilde{D}_{it}$ would be capturing a wealth-reelection interaction rather than a fiscal discretion effect. To assess this concern, I compute the pairwise correlations between $D_{it}$ and my measures of municipal wealth. The correlation between $D_{it}$ and GDP (PIB) per capita is 0.12, and between $D_{it}$ and the municipal HDI is 0.21. Both correlations are low, indicating that fiscal discretion captures substantial institutional variation in revenue composition that is not accounted for by municipal wealth. This variation reflects differences across municipalities in their participation in earmarked federal programs, local tax collection, and year-to-year changes in transfer rules, none of which are perfectly predicted by income or human development. We therefore conclude that $\tilde{D}_{it}$ is not simply a proxy for wealth, and that any results are not driven by confounding with municipal prosperity.

\subsubsection*{Heterogeneity in the Reelection Effect by Municipality Size}

I explore whether the reelection effect on discretionary spending composition varies systematically with municipality size. Smaller municipalities may have fewer electoral instruments available, such as less patronage capacity, smaller administrative payrolls, and less infrastructure budget. This could lead reelection-eligible mayors in fiscally constrained small cities to rely more heavily on direct transfers such as social assistance as their primary electoral tool. To test this hypothesis, I augment the baseline specification with an interaction term between reelection eligibility and log population:

\begin{equation}
    s^{d}_{ikt} = \beta_0
                + \beta_1 R_{it}
                + \beta_2 \tilde{D}_{it}
                + \beta_3 (R_{it} \times \tilde{D}_{it})
                + \beta_4 (R_{it} \times \ln P_i)
                + \beta_5 N_{it}
                + \beta_6 \ln P_i
                + \mathbf{X}_i'\boldsymbol{\gamma}
                + \mu_t
                + \varepsilon_{it}
\end{equation}

The coefficient $\beta_4$ on the size interaction is not statistically significant for any of the six spending categories examined. Moreover, its inclusion introduces multicollinearity between $R_{it}$, $\ln P_i$, and $R_{it} \times \ln P_i$ that inflates the standard errors on $\beta_1$, causing the main reelection effects for urbanismo and administration to lose statistical significance. I therefore conclude that there is no detectable heterogeneity in the reelection effect by municipality size and retain the updated baseline specification. The conditional clientelism results, which are driven by the interaction between reelection eligibility and fiscal discretion, are robust to this augmented specification, with $\beta_3$ remaining significant for social assistance ($p = 0.004$) and administration ($p = 0.012$) even in the presence of the size interaction.

\header{Results}

\subsection*{Reelection eligibility and spending composition}

\begin{table}[htbp]
\centering
\caption{\small{Effect of Reelection Eligibility on Discretionary Spending Shares ($\hat{\beta}_1$)}}
\label{tab:beta1}
\resizebox{\textwidth}{!}{%
\begin{tabular}{lcccccc}
\hline\hline
 & Health & Education & Urbanism & Social Assist. & Administration & Transport \\
\hline
$R_{it}$ & $-0.024$ & $0.121$ & $-0.944^{***}$ & $0.055$ & $0.886^{**}$ & $-0.018$ \\
 & $(0.286)$ & $(0.203)$ & $(0.298)$ & $(0.098)$ & $(0.367)$ & $(0.195)$ \\
\hline\hline
\multicolumn{7}{l}{\footnotesize Standard errors clustered at municipality level in parentheses.} \\
\multicolumn{7}{l}{\footnotesize $^{***}p<0.01$, $^{**}p<0.05$, $^{*}p<0.10$.} \\
\multicolumn{7}{l}{\footnotesize $\beta_1$ evaluated at sample mean of fiscal discretion ($\bar{D} = 71.4\%$).} \\
\multicolumn{7}{l}{\footnotesize All regressions include year fixed effects and municipality-level controls.} \\
\end{tabular}%
}
\end{table}


Table 1 presents the estimates of $\beta_1$, which is the effect of reelection eligibility on discretionary spending shares at the sample mean of fiscal discretion, across all six spending categories in this analysis. Reelection eligibility has a statistically and economically significant effect on two categories: urbanism and administration.

For urbanism, the estimated coefficient is $\hat{\beta}_1 = -0.944$ ($p = 0.002$), indicating that reelection-eligible mayors allocate approximately 0.94 percentage points less of their discretionary budget to urban infrastructure relative to term-limited mayors, holding all other controls constant. The negative sign is consistent with an electoral accountability mechanism: reelection-eligible mayors restrain infrastructure spending because corruption in public works --- a category historically associated with overbilling, kickbacks, and diversion of public funds in Brazilian municipalities --- carries electoral risk before the next election. Term-limited mayors, by contrast, face no such constraint and increase discretionary urbanism spending, consistent with rent extraction through inflated construction contracts when accountability disappears \cite{ferraz:2011}.

For administration, the estimated coefficient is $\hat{\beta}_1 = 0.886$ ($p = 0.016$), indicating that reelection-eligible mayors allocate approximately 0.89 percentage points more of their discretionary budget to general government overhead relative to term-limited mayors. I interpret this as evidence of political patronage: reelection-seeking mayors strategically expand the municipal payroll to generate political and electoral loyalty, consistent with the well-documented role of public employment as a tool of electoral mobilization in Brazilian municipalities \cite{colonnelli:2020}. This effect is identified entirely within the discretionary budget.

The two findings together \textit{suggest} a reallocation mechanism: reelection-eligible mayors shift discretionary resources away from urbanism and toward administration. This is consistent with a rational electoral strategy in which incumbents substitute a category prone to visible corruption (infrastructure) for a less visible but electorally productive form of spending (patronage hiring) when reelection is possible

Health, education, and transport show no significant reelection effect ($p > 0.50$ in all cases). This is consistent with constitutional earmarking constraining mayoral discretion in these categories, even within the discretionary budget: FUNDEB mandates minimum education spending, and SUS transfers create institutional pressures to maintain health spending levels, limiting the scope for politically motivated reallocation regardless of electoral incentives. I return to this point when discussing the impact of earmarked funds.

\subsection*{Fiscal discretion and spending composition}

\begin{table}[htbp]
\centering
\caption{Effect of Fiscal Discretion on Discretionary Spending Shares ($\hat{\beta}_2$)}
\label{tab:beta2}
\resizebox{\textwidth}{!}{%
\begin{tabular}{lcccccc}
\hline\hline
 & Health & Education & Urbanism & Social Assist. & Administration & Transport \\
\hline
$\tilde{D}_{it}$ & $-0.086^{**}$ & $0.064$ & $-0.106^{**}$ & $0.031^{**}$ & $0.027$ & $0.031$ \\
 & $(0.042)$ & $(0.035)$ & $(0.041)$ & $(0.015)$ & $(0.050)$ & $(0.026)$ \\
\hline\hline
\multicolumn{7}{l}{\footnotesize Standard errors clustered at municipality level in parentheses.} \\
\multicolumn{7}{l}{\footnotesize $^{***}p<0.01$, $^{**}p<0.05$, $^{*}p<0.10$.} \\
\multicolumn{7}{l}{\footnotesize $\tilde{D}_{it}$ is the demeaned discretionary share of total expenditure (sample mean = 71.4\%).} \\
\multicolumn{7}{l}{\footnotesize All regressions include year fixed effects and municipality-level controls.} \\
\end{tabular}%
}
\end{table}

Table 2 presents the estimates of $\beta_2$, which represents the effect of fiscal discretion on spending shares, measured relative to the sample mean of $\tilde{D}_{it} = 71.4\%$. Fiscal discretion has a statistically significant independent effect on three categories.

For social assistance, $\hat{\beta}_2 = 0.031$ ($p = 0.035$), indicating that municipalities with above-average fiscal discretion allocate proportionally more of their discretionary budget to welfare programs. For health, $\hat{\beta}_2 = -0.086$ ($p = 0.041$), and for urbanism, $\hat{\beta}_2 = -0.106$ ($p = 0.011$), both negative, indicating that municipalities with above-average fiscal discretion allocate proportionally less to these categories.

The negative coefficients for health and urbanism are interpretable as follows. When the discretionary share of the budget is above average, mayors have more room to fund other categories, and the relative importance of health and infrastructure in the discretionary budget falls. The positive coefficient for social assistance suggests that welfare programs are genuinely discretionary spending priorities that become more salient when mayors have fiscal room to allocate.

\subsection*{Conditional clientelism: the moderation effect of fiscal discretion}

\begin{table}[htbp]
\centering
\caption{Moderation Effect of Fiscal Discretion on Reelection Incentives ($\hat{\beta}_3$)}
\label{tab:beta3}
\resizebox{\textwidth}{!}{%
\begin{tabular}{lcccccc}
\hline\hline
 & Health & Education & Urbanism & Social Assist. & Administration & Transport \\
\hline
$R_{it} \times \tilde{D}_{it}$ & $0.070$ & $-0.066^{**}$ & $0.049$ & $-0.037^{***}$ & $-0.133^{**}$ & $-0.034$ \\
 & $(0.050)$ & $(0.034)$ & $(0.039)$ & $(0.014)$ & $(0.058)$ & $(0.024)$ \\
\hline\hline
\multicolumn{7}{l}{\footnotesize Standard errors clustered at municipality level in parentheses.} \\
\multicolumn{7}{l}{\footnotesize $^{***}p<0.01$, $^{**}p<0.05$, $^{*}p<0.10$.} \\
\multicolumn{7}{l}{\footnotesize A negative $\hat{\beta}_3$ indicates the reelection effect narrows for above-average discretion municipalities.} \\
\multicolumn{7}{l}{\footnotesize All regressions include year fixed effects and municipality-level controls.} \\
\end{tabular}%
}
\end{table}

Table 3 presents the estimates of $\beta_3$, which is the moderation effect of fiscal discretion on reelection incentives. The interaction term $R_{it} \times \tilde{D}_{it}$ is statistically significant for three categories: social assistance, administration, and education.

For social assistance, $\hat{\beta}_3 = -0.037$ ($p = 0.008$). For administration, $\hat{\beta}_3 = -0.133$ ($p = 0.022$). For education, $\hat{\beta}_3 = -0.066$ ($p = 0.049$). The negative sign in both cases means that the spending gap between reelection-eligible and term-limited mayors narrows for municipalities with above-average fiscal discretion and widens for municipalities with below-average fiscal discretion. This can be understood as mayors who can be reelected spend more on social assistance and administration when they have below-average fiscal discretion. Since $\beta_1$ is evaluated at the sample mean of $\tilde{D}_{it}$, the total reelection effect at any given level of fiscal discretion is:

\begin{equation}
    \frac{\partial s^{d}_{ikt}}{\partial R_{it}} = \hat{\beta}_1 + \hat{\beta}_3 \tilde{D}_{it}
\end{equation}

\begin{table}[htbp]
\centering
\caption{Total Reelection Effect by Level of Fiscal Discretion: $\hat{\beta}_1 + \hat{\beta}_3 \tilde{D}_{it}$}
\label{tab:conditional}
\begin{tabular}{lccc}
\hline\hline
 & $-1$ SD & Mean & $+1$ SD \\
 & ($\tilde{D}_{it} = -7.49$) & ($\tilde{D}_{it} = 0$) & ($\tilde{D}_{it} = +7.49$) \\
\hline
Education         & $+0.61$ & $+0.12$ & $-0.37$ \\
Social Assistance & $+0.33$ & $+0.05$ & $-0.22$ \\
Administration    & $+2.13$ & $+0.89$ & $-0.11$ \\
\hline\hline
\multicolumn{4}{l}{\footnotesize All values in percentage points of discretionary budget.} \\
\multicolumn{4}{l}{\footnotesize SD of $\tilde{D}_{it}$ = 7.49 percentage points.} \\
\multicolumn{4}{l}{\footnotesize Computed as $\hat{\beta}_1 + \hat{\beta}_3 \tilde{D}_{it}$ at each value of $\tilde{D}_{it}$.} \\
\end{tabular}
\end{table}

Table 4 evaluates this total effect at one standard deviation below and above the sample mean ($\tilde{D}_{it} = \pm 7.49$ percentage points). For administration, the total reelection effect ranges from $+2.13$ pp at one standard deviation below the mean to nearly zero ($-0.11$ pp) at one standard deviation above it. For social assistance, the total reelection effect ranges from $+0.33$ pp at one standard deviation below the mean to $-0.22$ pp above it. For education, it varies from $+0.61$ pp at one standard deviation below the mean to $-0.37$ above it.

The reelection effect on administration is more significant when compared to those of education and social assistance. Even though the percentage changes to social assistance and education are small, they are still statistically significant. However, the main reelection effect $\hat{\beta}_1$ is not significant for these two categories. This pattern --- a significant interaction without a significant main effect --- suggests that the reelection effect on education spending is heterogeneous across levels of fiscal discretion but averages out to zero at the sample mean. Specifically, reelection-eligible mayors in municipalities with below-average fiscal discretion allocate more of their discretionary budget to education and social assistance, while those with above-average fiscal discretion allocate less.  However, given that this result is only marginally significant and the main effect is not, it should be interpreted with caution and treated as suggestive rather than conclusive.

The results, mainly the ones for social assistance and administration, support a conditional clientelism interpretation: the form of electoral manipulation depends on the degree of fiscal constraint facing each mayor. Mayors with below-average fiscal discretion rely simultaneously on two electoral instruments: patronage hiring and social assistance transfers. It is important to note that, even though social assistance $\hat{\beta_3}$ has a smaller coefficient, even a slight percentage change in the total budget toward this category could affect the lives of financially vulnerable citizens and entice them to support the current mayor in the next election. Both channels become stronger as fiscal discretion falls below the mean. Mayors with above-average fiscal discretion show attenuated effects in both categories, suggesting that when budget room is ample, mayors have access to a wider range of electoral instruments and do not need to rely on direct hiring and transfers to mobilize support.

\subsection*{Impact of Earmarked funds: Health and Education}

Reelection eligibility has no significant effect on discretionary spending shares for health ($\hat{\beta}_1 = -0.024$, $p = 0.934$) or education ($\hat{\beta}_1 = 0.121$, $p = 0.550$), but discretion is significant for health ($\hat{\beta}_2 = -0.086$, $p = 0.041$) and the interaction term is significant for education ($\hat{\beta}_3 = -0.066$, $p = 0.049$). These are mixed results and should be interpreted with caution. The fact that terms related to discretion showed to be significant in both cases could suggest that the earmarked funds are valuable: a mayor with more power over the municipality's budget tends to divert money from education and health to other categories, so the binding earmarked funds deny him the possibility to move more money. However, it could also be interpreted as follows: a mayor chooses to invest in other categories that receive less federal support when given the choice, since education and health have already been accounted for.

\subsection*{Municipality Size and Controls}

Log population enters significantly and negatively for social assistance ($\hat{\beta}_6 = -0.553$, $p < 0.001$), administration ($\hat{\beta}_6 = -1.368$, $p < 0.001$), and transport ($\hat{\beta}_6 = -0.999$, $p < 0.001$). Larger municipalities allocate proportionally less of their discretionary budget to these categories, consistent with economies of scale in service delivery and the greater administrative complexity of larger cities, which may require a different spending mix than smaller municipalities. It also supports that smaller municipalities are more prone to clientelism and patronage, as some literature suggests \cite{toral:2024} \cite{nicther:2018}

Political competition, measured by the number of candidates in the relevant election ($N_{it}$), has a significant negative effect on health spending ($\hat{\beta}_5 = -0.336$, $p < 0.001$), could suggest that more competitive electoral environments are associated with less discretionary spending on health, possibly because mayors in competitive municipalities face stronger pressure to allocate resources across a wider range of visible categories rather than concentrating on any single one. However, since no other category presented a statistically significant coefficient for $\hat{\beta}_5$, the results suggest that this is an outlier, as we would expect that greater political pressure would have a higher impact on all categories. Also, the magnitude of this coefficient does not suggest any real/relevant interpretability (absolute value less than 1).

Among the socioeconomic controls, IDHM has the expected signs: positive and significant for education ($\hat{\beta} = 27.24$, $p < 0.001$) and negative for health ($\hat{\beta} = -27.76$, $p = 0.002$) and transport ($\hat{\beta} = -20.06$, $p < 0.001$). Wealthier and more developed municipalities spend proportionally more on education and less on health and transport in their discretionary budgets, consistent with differences in baseline service provision across the development spectrum.

\header{Conclusion}

This paper examines whether mayors who can be reelected spend differently than those who are term-bound in municipalities of São Paulo, while also identifying categories that mayors exploit to gain political support and sway votes. Using a panel of 643 municipalities in São Paulo state over the period 2017--2024, and exploiting Brazil's two-term mayoral limit as a source of exogenous variation in electoral incentives, I estimate a model of discretionary spending shares that separately identifies the effect of reelection eligibility, fiscal discretion, and their interaction.

My results yield two main findings. First, reelection eligibility significantly affects spending composition in two categories. Reelection-eligible mayors allocate approximately 0.94 percentage points less of their discretionary budget to urbanism relative to term-limited mayors, consistent with electoral accountability constraining opportunistic infrastructure spending: when reelection is possible, mayors refrain from construction contracts and kickbacks that term-limited mayors can exploit with relative electoral impunity. At the same time, reelection-eligible mayors allocate approximately 0.89 percentage points more to administration, consistent with political patronage: incumbents seeking reelection strategically expand the municipal payroll to generate electoral loyalty \cite{colonnelli:2020}. Together, these findings suggest that electoral accountability disciplines one form of opportunistic behavior --- rent extraction through public works --- while simultaneously encouraging another --- patronage hiring.

Second, I document a conditional clientelism mechanism through which fiscal discretion moderates the reelection effect on spending composition. Mayors with below-average fiscal discretion rely on two electoral instruments: patronage hiring and social assistance transfers. The total reelection effect on administration spending reaches $+2.13$ percentage points at one standard deviation below the mean of fiscal discretion, nearly three times the average effect, and falls to essentially zero for municipalities with above-average fiscal discretion. A similar pattern holds for social assistance. This finding, together with coefficients on log-population, is consistent with the broader clientelism literature documenting that patronage and vote-buying are more prevalent in smaller, fiscally constrained municipalities where personal networks dominate the allocation of public resources \cite{nichter:2018}.

These results have several implications for the design of fiscal institutions and electoral rules in decentralized settings. The finding that term limits generate opportunistic infrastructure spending --- rather than disciplining behavior as intended --- suggests that removing reelection incentives may have unintended consequences for fiscal governance. Some often advocate that term limits act as a check on incumbency advantage and entrenchment, but my results indicate that they create a window of reduced accountability in which mayors shift toward rent-extracting spending categories. This is consistent with the theoretical prediction of \citeasnoun{besley:1995} that lame-duck politicians maximize short-run payoff. More specifically, in the Brazilian context, they utilize infrastructure spending as rent extraction, a finding that adds to the empirical record on how accountability shapes the composition rather than aggregate expenditure. However, my results also suggest that incumbency advantage is present in São Paulo municipalities, since there is evidence that funds are being used by mayors who can be reelected differently than those who are term-bound, in accordance with the patronage and clientelism literature.

The conditional clientelism finding also has policy and organizational implications. If the instruments of electoral manipulation depend on fiscal discretion, then institutional reforms that reduce earmarking --- intended to give mayors more flexibility to respond to local needs --- may inadvertently expand the scope for politically motivated spending decisions, as my results suggest.

Several limitations of this analysis deserve mention. My identification strategy relies on the assumption that reelection eligibility is uncorrelated with unobserved determinants of spending composition, conditional on controls and year fixed effects. While Brazil's term limit rule is constitutionally mandated and determined by prior election outcomes, I cannot fully rule out that municipalities with first-term and second-term mayors differ in unobservable ways that affect spending decisions.

Future research could extend this analysis in 3 main directions. One could examine whether mayors from municipalities in São Paulo engage in greater spending in the years prior to election in a certain group of categories, indicating differences in spending patterns when the election approaches, which would likely not happen if mayors always acted in the public interest. Another possible extension would be utilizing Large Language Models (LLMs) to systematically analyze and categorize expenses according to relevant groups. The dataset contains more detailed descriptions of each expense that contain information on what/how the money was spent, which would allow for finer granularity and more interesting findings. Finally, this framework could be utilized to investigate spending patterns in municipalities in other states of Brazil and other regions in the world, looking for more datasets and a larger sample, which would help ground these empirical findings with even more statistical significance and broader applicability.

\bibliographystyle{econometrica}
\bibliography{Lucas_Sernik_final_draft}

\newpage

\header{Appendix}

\begin{table}[htbp]
\centering
\caption{Descriptive Statistics}
\label{tab:descriptive}
\resizebox{\textwidth}{!}{%
\begin{tabular}{lrrrrrrrr}
\hline\hline
Variable & $N$ & Mean & SD & Min & P25 & Median & P75 & Max \\
\hline
\multicolumn{9}{l}{\textit{Panel A: Dependent variables (discretionary spending share, \%)}} \\
\quad Health (\textit{Saúde}) & 5,144 & 26.348 & 5.175 & 0.00e+00 & 23.064 & 26.221 & 29.645 & 50.795 \\
\quad Education (\textit{Educação}) & 5,144 & 15.122 & 3.969 & 5.491 & 12.392 & 14.558 & 17.239 & 52.599 \\
\quad Urbanism & 5,116 & 11.746 & 5.121 & 0.00e+00 & 8.268 & 11.253 & 14.913 & 34.400 \\
\quad Social Assistance & 5,144 & 4.676 & 1.913 & 0.00e+00 & 3.302 & 4.469 & 5.683 & 22.365 \\
\quad Administration & 5,144 & 17.684 & 6.636 & 4.346 & 13.229 & 16.588 & 21.126 & 50.668 \\
\quad Transport & 3,726 & 3.287 & 3.102 & 0.00e+00 & 0.9832 & 2.454 & 4.514 & 20.927 \\
\hline
\multicolumn{9}{l}{\textit{Panel B: Key regressors}} \\
\quad Reelection eligibility ($R_{it}$) & 5,144 & 0.7970 & 0.4022 & 0.00e+00 & 1.000 & 1.000 & 1.000 & 1.000 \\
\quad Discretionary share ($D_{it}$, \%) & 5,144 & 71.443 & 7.489 & 38.571 & 66.591 & 71.697 & 76.419 & 94.398 \\
\quad Demeaned discretionary share ($\tilde{D}_{it}$) & 5,144 & -8.84e-15 & 7.489 & -32.873 & -4.852 & 0.2537 & 4.976 & 22.954 \\
\quad Interaction ($R_{it} \times \tilde{D}_{it}$) & 5,144 & -0.1034 & 6.764 & -32.873 & -3.528 & 0.00e+00 & 3.604 & 22.954 \\
\hline
\multicolumn{9}{l}{\textit{Panel C: Controls}} \\
\quad Log population ($\ln P_i$) & 5,144 & 9.632 & 1.382 & 6.691 & 8.547 & 9.452 & 10.543 & 14.016 \\
\quad Number of candidates ($N_{it}$) & 5,144 & 3.932 & 2.345 & 1.000 & 2.000 & 3.000 & 5.000 & 17.000 \\
\quad PIB per capita (R\$) & 5,144 & 1,762,171 & 5,498,120 & 18,877 & 96,738 & 246,197 & 1,024,226 & 54,363,858 \\
\quad IDHM & 5,144 & 0.7395 & 0.0324 & 0.6390 & 0.7190 & 0.7380 & 0.7610 & 0.8620 \\
\quad Share of children (\%) & 5,144 & 29.900 & 3.549 & 10.590 & 27.500 & 29.960 & 32.270 & 41.390 \\
\hline\hline
\multicolumn{9}{l}{\footnotesize Panel A variables are measured as percentage shares of total discretionary spending.} \\
\multicolumn{9}{l}{\footnotesize $D_{it}$ is the share of total expenditure funded by discretionary treasury resources (\%).} \\
\multicolumn{9}{l}{\footnotesize $\tilde{D}_{it} = D_{it} - \bar{D}$ where $\bar{D} = 71.4\%$ is the sample mean.} \\
\multicolumn{9}{l}{\footnotesize PIB per capita in 2010 BRL. Population from 2010 IBGE Census.} \\
\end{tabular}%
}
\end{table}
\begin{table}[htbp]
\centering
\caption{Full Regression Results: Discretionary Spending Shares by Category}
\label{tab:full}
\resizebox{\textwidth}{!}{%
\begin{tabular}{lcccccc}
\hline\hline
 & Health & Education & Urbanism & Social Assist. & Administration & Transport \\
\hline
$R_{it}$ & $-0.024$ & $0.121$ & $-0.944^{***}$ & $0.055$ & $0.886^{**}$ & $-0.018$ \\
 & $(0.286)$ & $(0.203)$ & $(0.298)$ & $(0.098)$ & $(0.367)$ & $(0.195)$ \\
$\tilde{D}_{it}$ & $-0.086^{**}$ & $0.064$ & $-0.106^{**}$ & $0.031^{**}$ & $0.027$ & $0.031$ \\
 & $(0.042)$ & $(0.035)$ & $(0.041)$ & $(0.015)$ & $(0.050)$ & $(0.026)$ \\
$R_{it} \times \tilde{D}_{it}$ & $0.070$ & $-0.066^{**}$ & $0.049$ & $-0.037^{***}$ & $-0.133^{**}$ & $-0.034$ \\
 & $(0.050)$ & $(0.034)$ & $(0.039)$ & $(0.014)$ & $(0.058)$ & $(0.024)$ \\
$\ln(P_i)$ & $0.287$ & $-0.074$ & $-0.190$ & $-0.553^{***}$ & $-1.368^{***}$ & $-0.999^{***}$ \\
 & $(0.226)$ & $(0.161)$ & $(0.216)$ & $(0.077)$ & $(0.288)$ & $(0.141)$ \\
$N_{it}$ & $-0.336^{***}$ & $-0.033$ & $-0.001$ & $-0.027$ & $0.031$ & $0.077$ \\
 & $(0.101)$ & $(0.077)$ & $(0.078)$ & $(0.026)$ & $(0.121)$ & $(0.049)$ \\
PIB per capita & $-0.000$ & $0.000$ & $-0.000^{**}$ & $0.000$ & $0.000$ & $0.000^{***}$ \\
 & $(0.000)$ & $(0.000)$ & $(0.000)$ & $(0.000)$ & $(0.000)$ & $(0.000)$ \\
IDHM & $-27.755^{***}$ & $27.235^{***}$ & $5.431$ & $-5.545^{*}$ & $-8.033$ & $-20.058^{***}$ \\
 & $(9.118)$ & $(6.256)$ & $(8.120)$ & $(3.154)$ & $(10.984)$ & $(5.239)$ \\
Pct. children & $0.098^{*}$ & $0.149^{***}$ & $0.158^{***}$ & $0.021$ & $-0.110$ & $-0.034$ \\
 & $(0.057)$ & $(0.042)$ & $(0.061)$ & $(0.023)$ & $(0.076)$ & $(0.042)$ \\
Constant & $42.547^{***}$ & $-8.780^{*}$ & $5.707$ & $13.517^{***}$ & $39.208^{***}$ & $28.135^{***}$ \\
 & $(6.657)$ & $(4.852)$ & $(6.031)$ & $(2.384)$ & $(8.169)$ & $(3.854)$ \\
\hline
Observations & $5{,}144$ & $5{,}144$ & $5{,}116$ & $5{,}144$ & $5{,}144$ & $3{,}726$ \\
Municipalities & $643$ & $643$ & $641$ & $643$ & $643$ & $524$ \\
Year FE & Yes & Yes & Yes & Yes & Yes & Yes \\
Clustered SE & Yes & Yes & Yes & Yes & Yes & Yes \\
R-squared & $0.076$ & $0.047$ & $0.037$ & $0.213$ & $0.088$ & $0.217$ \\
\hline\hline
\multicolumn{7}{l}{\footnotesize Standard errors clustered at municipality level in parentheses.} \\
\multicolumn{7}{l}{\footnotesize $^{***}p<0.01$, $^{**}p<0.05$, $^{*}p<0.10$.} \\
\multicolumn{7}{l}{\footnotesize Dependent variable: discretionary spending in category $k$ as share of total discretionary spending (\%).} \\
\multicolumn{7}{l}{\footnotesize $\tilde{D}_{it}$: demeaned discretionary share of total expenditure. $\ln(P_i)$: log 2010 Census population.} \\
\multicolumn{7}{l}{\footnotesize $N_{it}$: number of candidates in the relevant election, varying at municipality-term level.} \\
\end{tabular}%
}
\end{table}
\end{document}